\documentclass{beamer}

%
% Common preamble for all three parts.
%

%\usepackage[english]{babel}
\usepackage[polutonikogreek, italian]{babel}
\usepackage{amsmath}
\usepackage{color}
\usepackage[cache=false]{minted}
\usepackage{hyperref}
\usepackage{multicol}
\usepackage{tabularx}
\usepackage{tikz}
\usepackage{shapepar}

% only inline todonotes work
\usepackage{xkeyval}
\usepackage[textsize=small]{todonotes}
\presetkeys{todonotes}{inline}{}

\usetikzlibrary{shapes,arrows,positioning,shadows}

% no nav buttons
\usenavigationsymbolstemplate{}

\newcommand{\bftt}[1]{\textbf{\texttt{#1}}}
\newcommand{\comment}[1]{{\color[HTML]{008080}\textit{\textbf{\texttt{#1}}}}}
\newcommand{\cmd}[1]{{\color[HTML]{008000}\bftt{#1}}}
\newcommand{\bs}{\char`\\}
\newcommand{\cmdbs}[1]{\cmd{\bs#1}}
\newcommand{\lcb}{\char '173}
\newcommand{\rcb}{\char '175}
\newcommand{\cmdbegin}[1]{\cmdbs{begin\lcb}\bftt{#1}\cmd{\rcb}}
\newcommand{\cmdend}[1]{\cmdbs{end\lcb}\bftt{#1}\cmd{\rcb}}

\newcommand{\wllogo}{\textbf{Overleaf}}

% this is where the example source files are loaded from
% do not include a trailing slash
\newcommand{\fileuri}{https://raw.github.com/mirtexxan/latex-course/master/it}

\newcommand{\wlserver}{https://www.overleaf.com}
\newcommand{\wlnewdoc}[1]{\wlserver/docs?snip\_uri=\fileuri/#1\&splash=none}

\def\tikzname{Ti\emph{k}Z}

% from http://tex.stackexchange.com/questions/5226/keyboard-font-for-latex
\newcommand*\keystroke[1]{%
  \tikz[baseline=(key.base)]
    \node[%
      draw,
      fill=white,
      drop shadow={shadow xshift=0.25ex,shadow yshift=-0.25ex,fill=black,opacity=0.75},
      rectangle,
      rounded corners=2pt,
      inner sep=1pt,
      line width=0.5pt,
      font=\scriptsize\sffamily
    ](key) {#1\strut}
  ;
}
\newcommand{\keystrokebftt}[1]{\keystroke{\bftt{#1}}}

% stolen from minted.dtx
\newenvironment{exampletwoup}
  {\VerbatimEnvironment
   \begin{VerbatimOut}{example.out}}
  {\end{VerbatimOut}
   \setlength{\parindent}{0pt}
   \fbox{\begin{tabular}{l|l}
   \begin{minipage}{0.55\linewidth}
     \inputminted[fontsize=\small,resetmargins]{latex}{example.out}
   \end{minipage} &
   \begin{minipage}{0.35\linewidth}
     \input{example.out}
   \end{minipage}
   \end{tabular}}}

\newenvironment{exampletwouptiny}
  {\VerbatimEnvironment
   \begin{VerbatimOut}{example.out}}
  {\end{VerbatimOut}
   \setlength{\parindent}{0pt}
   \fbox{\begin{tabular}{l|l}
   \begin{minipage}{0.55\linewidth}
     \inputminted[fontsize=\scriptsize,resetmargins]{latex}{example.out}
   \end{minipage} &
   \begin{minipage}{0.35\linewidth}
     \setlength{\parskip}{6pt plus 1pt minus 1pt}%
     \raggedright\scriptsize\input{example.out}
   \end{minipage}
   \end{tabular}}}

\newenvironment{exampletwouptinynoframe}
  {\VerbatimEnvironment
   \begin{VerbatimOut}{example.out}}
  {\end{VerbatimOut}
   \setlength{\parindent}{0pt}
   \begin{tabular}{l|l}
   \begin{minipage}{0.55\linewidth}
     \inputminted[fontsize=\scriptsize,resetmargins]{latex}{example.out}
   \end{minipage} &
   \begin{minipage}{0.35\linewidth}
     \setlength{\parskip}{6pt plus 1pt minus 1pt}%
     \raggedright\scriptsize\input{example.out}
   \end{minipage}
   \end{tabular}}

\title{Muovere i primi passi con \LaTeX}
\author{Mirto Musci, PhD}
\institute{Assegnista di ricerca, Universit\`a di Pavia\\
Dipartimento di Ingegneria Industriale e dell'Informazione}
\titlegraphic{%
\includegraphics[height=1.8cm]{Unipv-logo}\hspace{1cm}
\includegraphics[height=1.8cm]{nuovo}%\hspace{1.2cm}
%\includegraphics[height=36pt]{overleaf}\\[1em]
}


\subtitle{Parte 1: Le basi}

\begin{document}

%%%%%%%%%%%%%%%%%%%%%%%%%%%%%%%%%%%%%%%%%%%%%%%%%%%%%%%%%%%%%%%%%%%%%%%%%%%%%%%
%%%%%%%%%%%%%%%%%%%%%%%%%%%%%%%%%%%%%%%%%%%%%%%%%%%%%%%%%%%%%%%%%%%%%%%%%%%%%%%
%%%%%%%%%%%%%%%%%%%%%%%%%%%%%%%%%%%%%%%%%%%%%%%%%%%%%%%%%%%%%%%%%%%%%%%%%%%%%%%
\begin{frame}
\titlepage
\end{frame}

%%%%%%%%%%%%%%%%%%%%%%%%%%%%%%%%%%%%%%%%%%%%%%%%%%%%%%%%%%%%%%%%%%%%%%%%%%%%%%%
%%%%%%%%%%%%%%%%%%%%%%%%%%%%%%%%%%%%%%%%%%%%%%%%%%%%%%%%%%%%%%%%%%%%%%%%%%%%%%%
%%%%%%%%%%%%%%%%%%%%%%%%%%%%%%%%%%%%%%%%%%%%%%%%%%%%%%%%%%%%%%%%%%%%%%%%%%%%%%%
\section{Introduzione}
\begin{frame}{Perch\'e \LaTeX{}?}
\begin{itemize}
\item Permette di realizzare bellissimi documenti professionali
\item Specialmente se contengono matematica
%
\item \`E stato creato da scienziati, per scienziati
\item \ldots ma si pu\`o usare in ogni campo del sapere!
%
\item \`E potentissimo --- ed estendibile a piacimento
\item Utilissimo per scrivere \structure{articoli, libri, tesi di laurea}, \ldots 
\end{itemize}
\end{frame}

%%%%%%%%%%%%%%%%%%%%%%%%%%%%%%%%%%%%%%%%%%%%%%%%%%%%%%%%%%%%%%%%%%%%%%%%%%%%%%%
%%%%%%%%%%%%%%%%%%%%%%%%%%%%%%%%%%%%%%%%%%%%%%%%%%%%%%%%%%%%%%%%%%%%%%%%%%%%%%%
%%%%%%%%%%%%%%%%%%%%%%%%%%%%%%%%%%%%%%%%%%%%%%%%%%%%%%%%%%%%%%%%%%%%%%%%%%%%%%%
\begin{frame}{Esempi}
\begin{equation}
\sum_{n = 0}^\infty \frac{(-1)^{n}}{2n+1} = 1 - \frac{1}{3} + \frac{1}{5} - \frac{1}{7} + \cdots = \frac{\pi}{4}
\end{equation}
\begin{equation}
\int_a^b \! f(x) dx = F(b) - F(a)
\end{equation}
\begin{equation}
f'(a)=\lim_{h\to 0}\frac{f(a+h)-f(a)}{h}
\end{equation}
\begin{equation}
\textstyle r=|z|=\sqrt{x^2+y^2}
\end{equation}
\begin{equation}
\begin{bmatrix} 1 & 2\\ 3 & 4\\ \end{bmatrix} \begin{bmatrix} 0 & 1\\ 0 & 0\\ \end{bmatrix}= \begin{bmatrix} 0 & 1\\ 0 & 3\\ \end{bmatrix}
\end{equation}
\vskip 2ex
\scriptsize e se imparate a fare matematica con Latex\ldots \url{http://latex2png.com/} 
\end{frame}

\begin{frame}{Esempi}
\shapepar{\heartshape}
Erano i capei d’oro a l’aura sparsi che ’n mille dolci nodi gli avolgea, e ’l vago lume oltra misura ardea di quei begli occhi, ch’or ne son sì scarsi; e ’l viso di pietosi color’ farsi, non so se vero o falso, mi parea: i’ che l’esca amorosa al petto avea,
qual meraviglia se di sùbito arsi? Non era l’andar suo cosa mortale, ma d’angelica forma; e le parole sonavan altro, che pur voce humana.
Uno spirto celeste, un vivo sole fu quel ch’i' vidi: e se non fosse or tale, piagha per allentar d’arco non sana.
\end{frame}

\begin{frame}{Esempi}
\bigskip\bigskip\bigskip
\foreignlanguage{polutonikogreek}{
\begin{verse}Th| p'anta dido'ush| ka`i >apolambano'ush| f'usei <o pepaideum'enos ka`i a>id'hmwn l'egei; <<d`os, <`o j'eleis, >ap'olabe, <`o j'eleic>>.
L'egei d`e to~uto o>u katajrasun'omenos, >all`a peijarq~wn m'onon ka`i e>uno~wn a>ut~h|.
\end{verse}}
\begin{flushright}
- Marco Aurelio, \emph{Ricordi}
\end{flushright}
\end{frame}

%%%%%%%%%%%%%%%%%%%%%%%%%%%%%%%%%%%%%%%%%%%%%%%%%%%%%%%%%%%%%%%%%%%%%%%%%%%%%%%
%%%%%%%%%%%%%%%%%%%%%%%%%%%%%%%%%%%%%%%%%%%%%%%%%%%%%%%%%%%%%%%%%%%%%%%%%%%%%%%
%%%%%%%%%%%%%%%%%%%%%%%%%%%%%%%%%%%%%%%%%%%%%%%%%%%%%%%%%%%%%%%%%%%%%%%%%%%%%%%
\begin{frame}{Come si pronuncia?}

\begin{itemize}
\item \TeX{} \`e stato creato a Stanford da Donald E. Knuth
\item Il nome deriva dalla radice greca di parole come
\textbf{\Large $$\tau\epsilon\chi\nu\acute{\eta}$$}
che significa \emph{arte} o \emph{tecnica}.
\begin{itemize}
\item la pronuncia dovrebbe essere \emph{tech} (come il tedesco \emph{Bach})
\item in italiano solitamente si usa \emph{tek}
\end{itemize}
\item \LaTeX{} \`e un set di macro per \TeX creato da Leslie Lamport

\begin{itemize}
\item oramai nessuno usa pi\`u \TeX{} base
\item la pronuncia usuale italiana \`e \emph{latek}.
\end{itemize}
\end{itemize}
\end{frame}

%%%%%%%%%%%%%%%%%%%%%%%%%%%%%%%%%%%%%%%%%%%%%%%%%%%%%%%%%%%%%%%%%%%%%%%%%%%%%%%
%%%%%%%%%%%%%%%%%%%%%%%%%%%%%%%%%%%%%%%%%%%%%%%%%%%%%%%%%%%%%%%%%%%%%%%%%%%%%%%
%%%%%%%%%%%%%%%%%%%%%%%%%%%%%%%%%%%%%%%%%%%%%%%%%%%%%%%%%%%%%%%%%%%%%%%%%%%%%%%

\begin{frame}[fragile]{Come funziona?}
\begin{itemize}
\item Un documento \LaTeX{} \`e composto di \texttt{testo semplice} misto a
\cmd{comandi} che ne descrivono la struttura e il significato.
\item L'applicazione \texttt{latex} \emph{compila} il testo e i comandi per produrre un documento perfettamente formattato.
\item In altre parole: quello che l'utente vede sullo schermo \structure{non \`e} il documento formattato.
\begin{itemize}
\item L'opposto di Word insomma, pi\`u simile ad HTML\ldots
\item What You See is What You Get vs. What You Mean
\end{itemize}

\end{itemize}
\vskip 1ex
\begin{center}
\begin{minted}[frame=single]{latex}
La rana in Spagna \emph{gracida} in campagna.
\end{minted}
\vskip 1ex
\tikz\node[single arrow,fill=gray,font=\ttfamily\bfseries,%
  rotate=270,xshift=-1em]{latex};
\vskip 1ex
\fbox{La rana in Spagna \emph{gracida} in campagna.}
\end{center}
\end{frame}

%%%%%%%%%%%%%%%%%%%%%%%%%%%%%%%%%%%%%%%%%%%%%%%%%%%%%%%%%%%%%%%%%%%%%%%%%%%%%%%
%%%%%%%%%%%%%%%%%%%%%%%%%%%%%%%%%%%%%%%%%%%%%%%%%%%%%%%%%%%%%%%%%%%%%%%%%%%%%%%
%%%%%%%%%%%%%%%%%%%%%%%%%%%%%%%%%%%%%%%%%%%%%%%%%%%%%%%%%%%%%%%%%%%%%%%%%%%%%%%
\begin{frame}[fragile]{Alcuni esempi di comandi}
\begin{exampletwoup}
\begin{itemize}
\item T\'e
\item Latte
\item Biscotti
\end{itemize}
\end{exampletwoup}
\vskip 2ex
\begin{exampletwoup}
\begin{figure}
\includegraphics{pulcino}
\end{figure}
\end{exampletwoup}
\vskip 2ex
\begin{exampletwoup}
\begin{equation}
\alpha + \beta + 1
\end{equation}
\end{exampletwoup}

\tiny{Immagine tratta da \url{http://www.andy-roberts.net/writing/latex/importing_images}}
\end{frame}

%%%%%%%%%%%%%%%%%%%%%%%%%%%%%%%%%%%%%%%%%%%%%%%%%%%%%%%%%%%%%%%%%%%%%%%%%%%%%%%
%%%%%%%%%%%%%%%%%%%%%%%%%%%%%%%%%%%%%%%%%%%%%%%%%%%%%%%%%%%%%%%%%%%%%%%%%%%%%%%
%%%%%%%%%%%%%%%%%%%%%%%%%%%%%%%%%%%%%%%%%%%%%%%%%%%%%%%%%%%%%%%%%%%%%%%%%%%%%%%

\begin{frame}[fragile]{In sintesi}

\begin{block}{Pregi}
\begin{itemize}
	\item \LaTeX{} comprende stili per la composizione di articoli o documenti di livello professionale.
	\item La composizione di formule matematiche \`e di altissimo livello.
	\item L'utente deve concentrarsi solo sul \structure{contenuto} del documento e non sulla sua \structure{forma} finale.
	\item Si possono generare note a pi\`e di pagina, riferimenti incrociati,
	bibliografie e indici in modo automatico.
	\item Migliaia di persone usano \LaTeX: è facile trovare supporto;
	\item \`E gratis.
\end{itemize}
\end{block}


\begin{alertblock}{Un cambio di paradigma}
Usate i comandi per descrivere `cio che \`e' e non `ci\`o che appare': concentratevi sul contenuto e Lasciate fare a \LaTeX{} il suo lavoro!
\end{alertblock}
\end{frame}


%%%%%%%%%%%%%%%%%%%%%%%%%%%%%%%%%%%%%%%%%%%%%%%%%%%%%%%%%%%%%%%%%%%%%%%%%%%%%%%
%%%%%%%%%%%%%%%%%%%%%%%%%%%%%%%%%%%%%%%%%%%%%%%%%%%%%%%%%%%%%%%%%%%%%%%%%%%%%%%
%%%%%%%%%%%%%%%%%%%%%%%%%%%%%%%%%%%%%%%%%%%%%%%%%%%%%%%%%%%%%%%%%%%%%%%%%%%%%%%
%\subsection{Compila?}
%
%\begin{frame}[fragile]{\insertsubsection{} Un classico programma C}
%
%Facciamo un paragone con il signor Pippo che vuole scrivere un programma
%nel suo linguaggio preferito, C
%
%\bigskip{}
%\begin{enumerate}
%\item Pippo scrive un documento di testo che chiamer\`a \cmd{pippo.c}.
%\item \label{enu:compila}Pippo d\`a sul terminale il comando
%\begin{minted}[frame=single]{bash}
%gcc pippo.c
%\end{minted}
%\item Se la compilazione d\`a errori, Pippo rivede il programma, lo corregge
%e ritorna al passo \ref{enu:compila}.
%\item Se la compilazione ha successo, viene prodotto un file eseguibile,
%\cmd{a.out}.
%\item Pippo prova a vedere se il programma fa ci\`o che desidera.
%\end{enumerate}
%\end{frame}
%
%%%%%%%%%%%%%%%%%%%%%%%%%%%%%%%%%%%%%%%%%%%%%%%%%%%%%%%%%%%%%%%%%%%%%%%%%%%%%%%%
%%%%%%%%%%%%%%%%%%%%%%%%%%%%%%%%%%%%%%%%%%%%%%%%%%%%%%%%%%%%%%%%%%%%%%%%%%%%%%%%
%%%%%%%%%%%%%%%%%%%%%%%%%%%%%%%%%%%%%%%%%%%%%%%%%%%%%%%%%%%%%%%%%%%%%%%%%%%%%%%%
%\begin{frame}[fragile]{\insertsubsection{} Il motore \LaTeX}
%
%Ora Pippo vuole scrivere la documentazione per il suo programma.
%\bigskip
%\begin{enumerate}
%\item Pippo scrive un documento di testo che chiamer\`a \cmd{pippo.tex}.
%\item \label{enu:compilalatex}Pippo d\`a sul terminale il comando
%
%\begin{minted}[frame=single]{bash}
%latex pippo.tex
%\end{minted}
%
%\item Se la compilazione d\`a errori, Pippo rivede il documento, lo corregge
%e ritorna al passo \ref{enu:compilalatex}.
%\item Se la compilazione ha successo, viene prodotto un file \cmd{pippo.dvi}.
%\item Pippo chiama il visualizzatore con
%
%\begin{minted}[frame=single]{bash}
%xdvi pippo
%\end{minted}
%
%e controlla che non ci siano errori concettuali.
%\end{enumerate}
%\end{frame}
%
%%%%%%%%%%%%%%%%%%%%%%%%%%%%%%%%%%%%%%%%%%%%%%%%%%%%%%%%%%%%%%%%%%%%%%%%%%%%%%%%
%%%%%%%%%%%%%%%%%%%%%%%%%%%%%%%%%%%%%%%%%%%%%%%%%%%%%%%%%%%%%%%%%%%%%%%%%%%%%%%%
%%%%%%%%%%%%%%%%%%%%%%%%%%%%%%%%%%%%%%%%%%%%%%%%%%%%%%%%%%%%%%%%%%%%%%%%%%%%%%%%
%\begin{frame}[fragile]{\insertsubsection{} pdf\LaTeX}
%
%Ora Pippo vuole scrivere la documentazione per il suo programma in
%un formato pi\`u comune.
%\bigskip
%\begin{enumerate}
%\item Pippo ha gi\`a il documento di testo chiamato \cmd{pippo.tex}.
%\item Pippo d\`a sul terminale il comando
%\begin{minted}[frame=single]{bash}
%pdflatex pippo.tex
%\end{minted}
%\item Viene prodotto un file \cmd{pippo.pdf}.
%\end{enumerate}
%\bigskip
%Dato che il documento \cmd{pippo.tex} \`e lo stesso di prima,
%il documento finale \`e solo la resa in PDF di quello precedente.
%\end{frame}

%%%%%%%%%%%%%%%%%%%%%%%%%%%%%%%%%%%%%%%%%%%%%%%%%%%%%%%%%%%%%%%%%%%%%%%%%%%%%%%
%%%%%%%%%%%%%%%%%%%%%%%%%%%%%%%%%%%%%%%%%%%%%%%%%%%%%%%%%%%%%%%%%%%%%%%%%%%%%%%
%%%%%%%%%%%%%%%%%%%%%%%%%%%%%%%%%%%%%%%%%%%%%%%%%%%%%%%%%%%%%%%%%%%%%%%%%%%%%%%
\section{Le basi}

%%%%%%%%%%%%%%%%%%%%%%%%%%%%%%%%%%%%%%%%%%%%%%%%%%%%%%%%%%%%%%%%%%%%%%%%%%%%%%%
%%%%%%%%%%%%%%%%%%%%%%%%%%%%%%%%%%%%%%%%%%%%%%%%%%%%%%%%%%%%%%%%%%%%%%%%%%%%%%%
%%%%%%%%%%%%%%%%%%%%%%%%%%%%%%%%%%%%%%%%%%%%%%%%%%%%%%%%%%%%%%%%%%%%%%%%%%%%%%%

\begin{frame}{Intermezzo: installare \LaTeX{}}

\begin{block}{Cosa scaricare}

\begin{itemize}
\item GNU/Linux: \TeX{}Live come sorgente o come pacchetto;
\item Mac OS X: Mac\TeX{};
\item Microsoft Windows: Mik\TeX{} (ce ne sono anche altri)
\end{itemize}
\end{block}

\begin{itemize}
\item Ottenuta la distribuzione, il codice si scrive con un editor
	\begin{itemize}
	\item Un \structure{qualsiasi} blocco note serve allo scopo
	\end{itemize}
\item Esistono anche ambienti integrati
	\begin{itemize}
	\item Editor `classici' come \TeX{}Works o \TeX{}Maker
	\item Strumenti `simil-word' come Lyx
	\item \ldots o web application come \wllogo
	\end{itemize}
\item Come \emph{compilare} dipende da programma a programma...
\end{itemize}

\end{frame}

%%%%%%%%%%%%%%%%%%%%%%%%%%%%%%%%%%%%%%%%%%%%%%%%%%%%%%%%%%%%%%%%%%%%%%%%%%%%%%%
%%%%%%%%%%%%%%%%%%%%%%%%%%%%%%%%%%%%%%%%%%%%%%%%%%%%%%%%%%%%%%%%%%%%%%%%%%%%%%%
%%%%%%%%%%%%%%%%%%%%%%%%%%%%%%%%%%%%%%%%%%%%%%%%%%%%%%%%%%%%%%%%%%%%%%%%%%%%%%%
\subsection{Iniziamo\ldots}
\begin{frame}[fragile]{\insertsubsection}
\begin{itemize}
\item Il documento \LaTeX{} minimale:
\vskip 2ex
\inputminted[frame=single]{latex}{basics.tex}
\item Tutti i comandi iniziano con un \emph{backslash} \keystrokebftt{\bs}.
\item Ogni documento inizia con un comando \cmdbs{documentclass}.
\item L'\emph{argomento} tra parentesi graffe \keystrokebftt{\{} \keystrokebftt{\}} indica a \LaTeX{} che tipo di documento stiamo creano: un \bftt{article}.
\item Il simbolo di percento \keystrokebftt{\%} d\`a inizio ad un \emph{commento}
--- \LaTeX{} ignorer\`a il resto della riga.
\end{itemize}
\end{frame}

%%%%%%%%%%%%%%%%%%%%%%%%%%%%%%%%%%%%%%%%%%%%%%%%%%%%%%%%%%%%%%%%%%%%%%%%%%%%%%%
%%%%%%%%%%%%%%%%%%%%%%%%%%%%%%%%%%%%%%%%%%%%%%%%%%%%%%%%%%%%%%%%%%%%%%%%%%%%%%%
%%%%%%%%%%%%%%%%%%%%%%%%%%%%%%%%%%%%%%%%%%%%%%%%%%%%%%%%%%%%%%%%%%%%%%%%%%%%%%%

\begin{frame}[fragile]{\insertsubsection{} con \wllogo}
\begin{itemize}
\item Overleaf \`e un webapp per scrivere documenti in \LaTeX.
\item Compila un sorgente \LaTeX{} e mostra i risultati in automatico e in
tempo reale. Permette di lavorare in \structure{condivisione}.
\item Senza preoccuparsi di dettagli come l'installazione o l'editor.
\vskip 2ex
\begin{center}
\fbox{\href{\wlnewdoc{basics.tex}}{%
Clicca qui per aprire \bftt{basics.tex} con \wllogo{}}}
\\[1ex]\scriptsize{}
Per migliore compatibilit\`a, usate \href{http://www.google.com/chrome}{Chrome} o un  \href{http://www.mozilla.org/en-US/firefox/new/}{FireFox} recente.
\end{center}
\vskip 2ex
\item Nel resto del corso, provate ad eseguire gli esempi, copiandoli
direttamente su Overleaf.
\item \textbf{No, davvero, \`e il miglior modo di imparare!}
\end{itemize}
\end{frame}

%%%%%%%%%%%%%%%%%%%%%%%%%%%%%%%%%%%%%%%%%%%%%%%%%%%%%%%%%%%%%%%%%%%%%%%%%%%%%%%
%%%%%%%%%%%%%%%%%%%%%%%%%%%%%%%%%%%%%%%%%%%%%%%%%%%%%%%%%%%%%%%%%%%%%%%%%%%%%%%
%%%%%%%%%%%%%%%%%%%%%%%%%%%%%%%%%%%%%%%%%%%%%%%%%%%%%%%%%%%%%%%%%%%%%%%%%%%%%%%
\subsection{Inserire il testo}
\begin{frame}[fragile]{\insertsubsection{}}
\small
\begin{itemize}
\item Tutto il testo di un qualunque documento va inserito tra \cmdbegin{document}
e \cmdend{document}.
\item Nella maggior parte dei casi, potete inserire testo normalmente.
\begin{exampletwouptiny}
Le parole sono separate da uno
o pi\`u spazi.

I paragrafi sono separati da una
o pi\`u righe vuote.
\end{exampletwouptiny}
\item Lo spazio nel file sorgente viene \emph{aggregato} nell'output.
\begin{exampletwouptiny}
La   rana       in Spagna
gracida in		 montagna.
\end{exampletwouptiny}
\item \keystrokebftt{\bs}\keystrokebftt{\bs} forza l'andata a capo, ma non crea un nuovo paragrafo.
\end{itemize}
\end{frame}

%%%%%%%%%%%%%%%%%%%%%%%%%%%%%%%%%%%%%%%%%%%%%%%%%%%%%%%%%%%%%%%%%%%%%%%%%%%%%%%
%%%%%%%%%%%%%%%%%%%%%%%%%%%%%%%%%%%%%%%%%%%%%%%%%%%%%%%%%%%%%%%%%%%%%%%%%%%%%%%
%%%%%%%%%%%%%%%%%%%%%%%%%%%%%%%%%%%%%%%%%%%%%%%%%%%%%%%%%%%%%%%%%%%%%%%%%%%%%%%
\begin{frame}[fragile]{\insertsubsection{}: caratteri speciali}
\small
\begin{itemize}
\item Le virgolette richiedono attenzione: va usato un apostrofo rovesciato \keystroke{\`{}} a sinistra e un apostrofo semplice \keystroke{\'{}} a destra
\begin{exampletwouptiny}
Virgolette semplici: `testo'.

Virgolette doppie: ``testo''.
\end{exampletwouptiny}

\item Alcuni simboli comuni hanno significato speciale in \LaTeX:\\[1ex]
\begin{tabular}{cl|l}
\keystrokebftt{\%} & percento    & commenti \\
\keystrokebftt{\#} & cancelletto & comandi custom \\
\keystrokebftt{\&} & ampersand   & tabelle \\
\keystrokebftt{\$} & dollaro     & matematica \\
\end{tabular}
\item Se provate ad inserirli direttamente, otterrete un messaggio di errore.
Se volete mostrarli nel documento dovete fare \structure{\emph{escape}},
precedendoli con un backslash \keystrokebftt{\bs}\\
\begin{exampletwouptiny}
\$\%\&\#!
\end{exampletwouptiny}
\end{itemize}

\end{frame}

%%%%%%%%%%%%%%%%%%%%%%%%%%%%%%%%%%%%%%%%%%%%%%%%%%%%%%%%%%%%%%%%%%%%%%%%%%%%%%%
%%%%%%%%%%%%%%%%%%%%%%%%%%%%%%%%%%%%%%%%%%%%%%%%%%%%%%%%%%%%%%%%%%%%%%%%%%%%%%%
%%%%%%%%%%%%%%%%%%%%%%%%%%%%%%%%%%%%%%%%%%%%%%%%%%%%%%%%%%%%%%%%%%%%%%%%%%%%%%%

\begin{frame}[fragile]{\insertsubsection{}: accenti}
\small
\begin{itemize}
\item Il caratteri base di \LaTeX{} sono per la lingua inglese (ASCII), ma per molte altre lingue, gli accenti sono importanti.
\begin{itemize}
\item \cmd{\bs`} (backslash+backtick) si usa per l'accento grave. Sfortunatamente l'apostrofo rovesciato non \`e semplice da inserire con la tastiera italiana.
\item \cmd{\bs'} (backslash+apostrofo) si usa per l'accento acuto.
\end{itemize}
\item Oppure si possono inserire direttamente le lettere accentate usando:
\mintinline{latex}!\usepackage[utf8]{inputenc}!
\item Esempio:
\begin{exampletwouptiny}
R\`{e}n\'{e} Descartes \'{e} noto
alla latina come `Cartesio', 
fu scienziato e ``filosofo''.
Mor\`{i} a Stoccolma di polmonite.
\end{exampletwouptiny}

\item Altri glifi non presenti in italiano, e relativi comandi:
\begin{exampletwouptiny}
Dieresi: Fl\"ugel\\
Circonflesso: H\^opital\\
Tilde: Vamos a ga\~nar\\
\end{exampletwouptiny}
\end{itemize}
\end{frame}

%%%%%%%%%%%%%%%%%%%%%%%%%%%%%%%%%%%%%%%%%%%%%%%%%%%%%%%%%%%%%%%%%%%%%%%%%%%%%%%
%%%%%%%%%%%%%%%%%%%%%%%%%%%%%%%%%%%%%%%%%%%%%%%%%%%%%%%%%%%%%%%%%%%%%%%%%%%%%%%
%%%%%%%%%%%%%%%%%%%%%%%%%%%%%%%%%%%%%%%%%%%%%%%%%%%%%%%%%%%%%%%%%%%%%%%%%%%%%%%
\begin{frame}[fragile]{Interludio: comandi}
\begin{itemize}
\item Tutti i comandi iniziano con \keystrokebftt{\bs} seguito da un simbolo (come per gli accenti) o una parola. Attenti a typos e maiuscole!
\item Effetto solo sull'argomento o sull'intero `paragrafo'

\begin{exampletwouptiny}
Parola in \emph{italico}.\\
\it Paragrafo  in italico.
\end{exampletwouptiny}
\vskip 1ex
\item I comandi possono avere argomenti obbligatori e opzionali
\vskip 1ex
\begin{exampletwouptiny}
\framebox        {un esempio} sciocco

\framebox[2cm]   {un esempio} sciocco

\framebox[2cm][r]{un esempio} sciocco

\end{exampletwouptiny}
\vskip 1ex

\item Caso speciale: la tilde \keystrokebftt{\textasciitilde}. Comando di un carattere, serve per specificare uno spazio che non si pu\`o separare

\begin{exampletwouptiny}
In riferimento all'es. a pag. 2\\
In riferimento all'es. a pag.~2
\end{exampletwouptiny}
\end{itemize}
\end{frame}

%%%%%%%%%%%%%%%%%%%%%%%%%%%%%%%%%%%%%%%%%%%%%%%%%%%%%%%%%%%%%%%%%%%%%%%%%%%%%%%
%%%%%%%%%%%%%%%%%%%%%%%%%%%%%%%%%%%%%%%%%%%%%%%%%%%%%%%%%%%%%%%%%%%%%%%%%%%%%%%
%%%%%%%%%%%%%%%%%%%%%%%%%%%%%%%%%%%%%%%%%%%%%%%%%%%%%%%%%%%%%%%%%%%%%%%%%%%%%%%
\begin{frame}[fragile]{Gestire gli errori}
\begin{itemize}
\item \LaTeX{} pu\`o confondersi nel compilare un documento.
Se succede, si interrompe con un \structure{messaggio di errore}.
\item Dovete correggere gli errori, se volete avere qualche speranza di
produrre un documento.
\item Per esempio, se provate a scrivere \cmdbs{epmh} invece di \cmdbs{emph},
\LaTeX{} si lamenter\`a con un errore \alert{\texttt{undefined control sequence}} dato che ``epmh'' non esiste come comando.
\end{itemize}
\begin{block}{Qualche consiglio sugli errori}
\begin{enumerate}
\item Niente panico! Succede a tutti.
\item Correggeli immediatamente --- se quello che avete appena scritto ha causato
un errore, perlomeno sapete da dove partire per il debugging.
\item Se ci sono errori multipli, correggeteli uno alla volta iniziando dal primo 
--- potrebbero essere errori a cascata.
\end{enumerate}
\end{block}
\end{frame}

%%%%%%%%%%%%%%%%%%%%%%%%%%%%%%%%%%%%%%%%%%%%%%%%%%%%%%%%%%%%%%%%%%%%%%%%%%%%%%%
%%%%%%%%%%%%%%%%%%%%%%%%%%%%%%%%%%%%%%%%%%%%%%%%%%%%%%%%%%%%%%%%%%%%%%%%%%%%%%%
%%%%%%%%%%%%%%%%%%%%%%%%%%%%%%%%%%%%%%%%%%%%%%%%%%%%%%%%%%%%%%%%%%%%%%%%%%%%%%%
\begin{frame}[fragile]{Esercizio 1}

\begin{block}{Scrivi questo in \LaTeX{}:%
\footnote{\url{http://en.wikipedia.org/wiki/Economy_of_the_United_States}}}
Nel Marzo 2006, il Congresso aument\`o la soglia di
\$790 miliardi per un totale di \$8970 miliardi,
che rappresentava circa il 68\% del PIL. Il 4 Ottobre
2008, l'``Emergency Economic Stabilization
Act'' aument\`o ulteriormente il tetto del debito
a \$11300 miliardi.
\end{block}
\vskip 2ex
\begin{center}
\fbox{\href{\wlnewdoc{basics-exercise-1.tex}}{%
Clicca per aprire \bftt{basics-exercise-1.tex} in \wllogo{}}}
\end{center}

\begin{itemize}
\item Suggerimento: attenti ai caratteri speciali!
\item Non dimenticate virgolette e accenti \ldots
\item Cliccate,
\fbox{\href{\wlnewdoc{basics-exercise-1-solution.tex}}{%
qui}} per la soluzione (\bftt{basics-exercise-1-solution.tex}).
\end{itemize}
\end{frame}

%%%%%%%%%%%%%%%%%%%%%%%%%%%%%%%%%%%%%%%%%%%%%%%%%%%%%%%%%%%%%%%%%%%%%%%%%%%%%%%
%%%%%%%%%%%%%%%%%%%%%%%%%%%%%%%%%%%%%%%%%%%%%%%%%%%%%%%%%%%%%%%%%%%%%%%%%%%%%%%
%%%%%%%%%%%%%%%%%%%%%%%%%%%%%%%%%%%%%%%%%%%%%%%%%%%%%%%%%%%%%%%%%%%%%%%%%%%%%%%
\subsection{Ambienti matematici}
\begin{frame}[fragile]{\insertsubsection{}: il dollaro}
\begin{itemize}
\item Come mai il dollaro \keystrokebftt{\$} \`e un simbolo speciale?
Lo si usa per separare l'ambiente matematico dal testo.\\[1ex]
\begin{exampletwouptiny}
% senza ambiente matematico:
Siano a e b due interi positivi
diversi, e sia c = a - b + 1.

% molto meglio, no?
Siano $a$ e $b$ due interi positivi
diversi, e sia $c = a - b + 1$.
\end{exampletwouptiny}
\item I simboli dollaro vanno sempre usati in coppia --- un per aprire l'ambiente
matematico, l'altro per chiuderlo.
\item Come al solito, \LaTeX{} gestisce la spaziatura in automatico,
ignorando quella dell'utente.
\begin{exampletwouptiny}
Sia $y=mx+b$ dove \ldots

Sia $y = m x + b$ dove \ldots
\end{exampletwouptiny}
\end{itemize}
\end{frame}

%%%%%%%%%%%%%%%%%%%%%%%%%%%%%%%%%%%%%%%%%%%%%%%%%%%%%%%%%%%%%%%%%%%%%%%%%%%%%%%
%%%%%%%%%%%%%%%%%%%%%%%%%%%%%%%%%%%%%%%%%%%%%%%%%%%%%%%%%%%%%%%%%%%%%%%%%%%%%%%
%%%%%%%%%%%%%%%%%%%%%%%%%%%%%%%%%%%%%%%%%%%%%%%%%%%%%%%%%%%%%%%%%%%%%%%%%%%%%%%
\begin{frame}[fragile]{\insertsubsection{}: Notazione}
\begin{itemize}
\item Usate circonflesso \keystrokebftt{\^} per gli apici e \emph{underscore} \keystrokebftt{\_} per i pedici. Sono altri due comandi speciali di un solo carattere.

\begin{exampletwouptiny}
$y = c_2 x^2 + c_1 x + c_0$
\end{exampletwouptiny}
\vskip 2ex

\item Usate le parentesi graffe \keystrokebftt{\{} \keystrokebftt{\}} per apici e 
pedici pi\`u lunghi. Racchiudono gli argomenti dei comandi \keystrokebftt{\^} e \keystrokebftt{\_}.
\begin{exampletwouptiny}
$F_n = F_n-1 + F_n-2$     % oops!

$F_n = F_{n-1} + F_{n-2}$ % ok!
\end{exampletwouptiny}
\vskip 2ex

\item \LaTeX{} offre molti comandi per rappresentare le lettere greche e la notazione pi\`u comune. Ad esempio:

\begin{exampletwouptiny}
$\mu = \alpha \int_t e^{q/rt} dt$

$\Omega = \sum_{k=1}^{n} \omega_k$
\end{exampletwouptiny}
\end{itemize}
\end{frame}

%%%%%%%%%%%%%%%%%%%%%%%%%%%%%%%%%%%%%%%%%%%%%%%%%%%%%%%%%%%%%%%%%%%%%%%%%%%%%%%
%%%%%%%%%%%%%%%%%%%%%%%%%%%%%%%%%%%%%%%%%%%%%%%%%%%%%%%%%%%%%%%%%%%%%%%%%%%%%%%
%%%%%%%%%%%%%%%%%%%%%%%%%%%%%%%%%%%%%%%%%%%%%%%%%%%%%%%%%%%%%%%%%%%%%%%%%%%%%%%
\begin{frame}[fragile]{\insertsubsection{}: Equazioni non in linea}
\begin{itemize}
\item Un'equazione lunga e complessa, andrebbe mostrata a parte usando
\cmdbegin{equation} e \cmdend{equation}.
\begin{itemize}
\item In alternativa: \mintinline{latex}!$$...$$ o \[...\]!\\[2ex]
\end{itemize}
\begin{exampletwouptiny}
Le radici di un'equazione quadratica
sono date da
\begin{equation}
x = \frac{-b \pm \sqrt{b^2 - 4ac}}
         {2a}
\end{equation}
dove $a$, $b$ e $c$ sono \ldots
\end{exampletwouptiny}
\vskip 1em
{\scriptsize Attenzione: \LaTeX{} ignora gli spazi in ambiente matematico, ma non
\`e in grado di gestire le line vuote --- non usatele!}
\end{itemize}
\end{frame}

%%%%%%%%%%%%%%%%%%%%%%%%%%%%%%%%%%%%%%%%%%%%%%%%%%%%%%%%%%%%%%%%%%%%%%%%%%%%%%%
%%%%%%%%%%%%%%%%%%%%%%%%%%%%%%%%%%%%%%%%%%%%%%%%%%%%%%%%%%%%%%%%%%%%%%%%%%%%%%%
%%%%%%%%%%%%%%%%%%%%%%%%%%%%%%%%%%%%%%%%%%%%%%%%%%%%%%%%%%%%%%%%%%%%%%%%%%%%%%%
\begin{frame}[fragile]{Intermezzo: Ambienti}
\begin{itemize}
\item \bftt{equation} \`e un \emph{ambiente} --- un \structure{contesto sematico}.
\item Un comando pu\`o produrre risultati diversi in contesti diversi.
\begin{exampletwouptiny}
Possiamo scrivere
$ \Omega = \sum_{k=1}^{n} \omega_k $
in linea, o usare un ambiente
\begin{equation}
  \Omega = \sum_{k=1}^{n} \omega_k
\end{equation}
per mostrarlo.
\end{exampletwouptiny}
\vskip 2ex
\item Notate che $\Sigma$ \`e pi\`u grande nell'ambiente \bftt{equation}, e
che apici e pedici cambiano di posizione, nonostante si siano usati gli
stessi comandi
\vskip 1ex
{\scriptsize Di fatto \mintinline{latex}!$...$! equivale a \cmdbegin{math}\bftt{...}\cmdend{math} e quindi \`e un ambiente}
\end{itemize}
\end{frame}

%%%%%%%%%%%%%%%%%%%%%%%%%%%%%%%%%%%%%%%%%%%%%%%%%%%%%%%%%%%%%%%%%%%%%%%%%%%%%%%
%%%%%%%%%%%%%%%%%%%%%%%%%%%%%%%%%%%%%%%%%%%%%%%%%%%%%%%%%%%%%%%%%%%%%%%%%%%%%%%
%%%%%%%%%%%%%%%%%%%%%%%%%%%%%%%%%%%%%%%%%%%%%%%%%%%%%%%%%%%%%%%%%%%%%%%%%%%%%%%
\begin{frame}[fragile]{Intermezzo: Ambienti}
\begin{itemize}
\item I comandi \cmdbs{begin} e \cmdbs{end} si possono usare per creare molti
ambienti diversi.
\item Gli ambienti \bftt{itemize} ed \bftt{enumerate} generano liste.
\begin{exampletwouptiny}
\begin{itemize} % elenco puntato
\item Biscotti
\item T\'e
\end{itemize}

\begin{enumerate} % elenco numerato
\item Biscotti
\item T\'e
\end{enumerate}
\end{exampletwouptiny}
\vskip 1ex
\item E \bftt{quote} e \bftt{quotation} permettono di inserire citazioni\ldots
\begin{exampletwouptiny}
\begin{quotation}
Four score and seven years ago
our fathers brought [\ldots]
\end{quotation}
\end{exampletwouptiny}
\end{itemize}
\end{frame}

%%%%%%%%%%%%%%%%%%%%%%%%%%%%%%%%%%%%%%%%%%%%%%%%%%%%%%%%%%%%%%%%%%%%%%%%%%%%%%%
%%%%%%%%%%%%%%%%%%%%%%%%%%%%%%%%%%%%%%%%%%%%%%%%%%%%%%%%%%%%%%%%%%%%%%%%%%%%%%%
%%%%%%%%%%%%%%%%%%%%%%%%%%%%%%%%%%%%%%%%%%%%%%%%%%%%%%%%%%%%%%%%%%%%%%%%%%%%%%%
\begin{frame}[fragile]{Intermezzo: Pacchetti}

\begin{itemize}
\item Tutti i comandi e gli ambienti mostrati fino ad adesso, sono parte di
\LaTeX{} \emph{base}.

\item I \emph{pacchetti} sono librerie di comandi e ambienti aggiuntivi:
ci sono migliaia di pacchetti liberamente disponibili.

\item I pacchetti che vogliamo usare vanno caricati esplicitamente usando il comando
\cmdbs{usepackage} nel \structure{preambolo}, ovvero prima di \cmdbegin{document}, ma dopo \mintinline{latex}|\documentclass{...}|

\item Esempio: \bftt{amsmath} della American Mathematical Society.
\begin{minted}[fontsize=\small,frame=single]{latex}
\documentclass{article}
\usepackage{amsmath} % preambolo
\begin{document}
% da ora in poi possiamo usare i comandi di amsmath...
\end{document}
\end{minted}
\end{itemize}
\end{frame}

%%%%%%%%%%%%%%%%%%%%%%%%%%%%%%%%%%%%%%%%%%%%%%%%%%%%%%%%%%%%%%%%%%%%%%%%%%%%%%%
%%%%%%%%%%%%%%%%%%%%%%%%%%%%%%%%%%%%%%%%%%%%%%%%%%%%%%%%%%%%%%%%%%%%%%%%%%%%%%%
%%%%%%%%%%%%%%%%%%%%%%%%%%%%%%%%%%%%%%%%%%%%%%%%%%%%%%%%%%%%%%%%%%%%%%%%%%%%%%%
\begin{frame}[fragile]{\insertsubsection{}: Esempi con \bftt{amsmath}}
\begin{itemize}
\item Usa \bftt{equation*} per inserire equazioni non numerate.
\begin{exampletwouptiny}
\begin{equation*}
  \Omega = \sum_{k=1}^{n} \omega_k
\end{equation*}
\end{exampletwouptiny}
\item \LaTeX{} tratta lettere adiacenti come variabili moltiplicate tra di loro, ma non \`e sempre desiderabile: \bftt{amsmath} definisce comandi per la maggior parte delle funzioni matematiche.
\begin{exampletwouptiny}
\begin{equation*} % sbagliato!
 min_{x,y} (1-x)^2 + 100(y-x^2)^2
\end{equation*}
\begin{equation*} % giusto!
\min_{x,y}{(1-x)^2 + 100(y-x^2)^2}
\end{equation*}
\end{exampletwouptiny}
\item Per quelle non predefinite, si usa \cmdbs{operatorname}.
\begin{exampletwouptiny}
\begin{equation*}
\beta_i =
\frac{\operatorname{Cov}(R_i, R_m)}
     {\operatorname{Var}(R_m)}
\end{equation*}
\end{exampletwouptiny}
\end{itemize}
\end{frame}

%%%%%%%%%%%%%%%%%%%%%%%%%%%%%%%%%%%%%%%%%%%%%%%%%%%%%%%%%%%%%%%%%%%%%%%%%%%%%%%
%%%%%%%%%%%%%%%%%%%%%%%%%%%%%%%%%%%%%%%%%%%%%%%%%%%%%%%%%%%%%%%%%%%%%%%%%%%%%%%
%%%%%%%%%%%%%%%%%%%%%%%%%%%%%%%%%%%%%%%%%%%%%%%%%%%%%%%%%%%%%%%%%%%%%%%%%%%%%%%
\begin{frame}[fragile]{\insertsubsection{}: Esempi con \bftt{amsmath}}
\begin{itemize}{\small
\item Allinea una sequenza di equazioni con il simbolo di uguale
\begin{align*}
(x+1)^3 &= (x+1)(x+1)(x+1) \\
        &= (x+1)(x^2 + 2x + 1) \\
        &= x^3 + 3x^2 + 3x + 1
\end{align*}
utilizzando l'ambiente \bftt{align*}.

% for whatever reason, this doesn't play well with the twoup environment
\begin{minted}[fontsize=\small,frame=single]{latex}
\begin{align*}
(x+1)^3 &= (x+1)(x+1)(x+1) \\
        &= (x+1)(x^2 + 2x + 1) \\
        &= x^3 + 3x^2 + 3x + 1
\end{align*}
\end{minted}
\item Una ampersand \keystrokebftt{\&} separa la colonna sinistra (prima di
$=$) dalla colonna destra (dopo di $=$).
\item Un doppio backslash \keystrokebftt{\bs}\keystrokebftt{\bs} inizia una
nuova linea.
}\end{itemize}
\end{frame}


%%%%%%%%%%%%%%%%%%%%%%%%%%%%%%%%%%%%%%%%%%%%%%%%%%%%%%%%%%%%%%%%%%%%%%%%%%%%%%%
%%%%%%%%%%%%%%%%%%%%%%%%%%%%%%%%%%%%%%%%%%%%%%%%%%%%%%%%%%%%%%%%%%%%%%%%%%%%%%%
%%%%%%%%%%%%%%%%%%%%%%%%%%%%%%%%%%%%%%%%%%%%%%%%%%%%%%%%%%%%%%%%%%%%%%%%%%%%%%%
\begin{frame}[fragile]{Esercizio 2}

\begin{block}{Scrivi questo in \LaTeX:}
Siano $X_1, X_2, \ldots, X_n$ una serie di variabili casuali indipendenti ed
identicamente distribuite tali per cui $\operatorname{E}[X_i] = \mu$ e
$\operatorname{Var}[X_i] = \sigma^2 < \infty$, con media
\begin{equation*}
S_n = \frac{1}{n}\sum_{i}^{n} X_i
\end{equation*}
Per $n$ che tende ad infinito, le variabili casuali
$\sqrt{n}(S_n - \mu)$ convergono in senso distribuzionale ad una gaussiana $N(0, \sigma^2)$.
\end{block}
\vskip 2ex
\begin{center}
\fbox{\href{\wlnewdoc{basics-exercise-2.tex}}{%
Clicca per aprire \bftt{basics-exercise-2.tex} in \wllogo{}}}
\end{center}
\begin{itemize}
\item Suggerimento: il comando per $\infty$ \`e \cmdbs{infty}.
\item Qui potete trovare la 
\fbox{\href{\wlnewdoc{basics-exercise-2-solution.tex}}{%
soluzione}} (\bftt{basics-exercise-2-solution.tex}).
\end{itemize}
\end{frame}

%%%%%%%%%%%%%%%%%%%%%%%%%%%%%%%%%%%%%%%%%%%%%%%%%%%%%%%%%%%%%%%%%%%%%%%%%%%%%%%
%%%%%%%%%%%%%%%%%%%%%%%%%%%%%%%%%%%%%%%%%%%%%%%%%%%%%%%%%%%%%%%%%%%%%%%%%%%%%%%
%%%%%%%%%%%%%%%%%%%%%%%%%%%%%%%%%%%%%%%%%%%%%%%%%%%%%%%%%%%%%%%%%%%%%%%%%%%%%%%
\begin{frame}{Fine della prima parte}
\begin{itemize}
\item \structure{Congratulazioni!} Avete imparato a\ldots
\begin{itemize}
\item Inserire testo in \LaTeX.
\item Utilizzare i comandi di base.
\item Gestire gli errori via via che compaiono.
\item Scrivere della bellissima matematica.
\item Utilizzare alcuni ambienti.
\item Caricare pacchetti.
\end{itemize}
\item Non \`e fantastico?
\item Nella \structure{seconda parte}, impareremo ad usare \LaTeX{} per scrivere
documenti strutturati con sezioni, riferimenti incrociati, figure,
tabelle, bibliografia\ldots. Alla prossima!
\end{itemize}
\end{frame}

\end{document}
